\begin{corollary}[Stokes torsion 的一回路计数律：直和回路格的球截断]\label{cor:hypercube-cycle-lattice-directsum-counting}
令 \(G\) 为有限连通多重图，并沿用定理 \ref{thm:hypercube-cycle-lattice-theta-asymptotic} 的回路格 \(\Lambda(G)\subset \RR^{E(G)}\)、秩 \(r\) 与生成树个数 \(\tau(G)\) 的记号。
对整数 \(m\ge 1\) 定义直和格
\[
\Lambda(G)^{\oplus m}:=\underbrace{\Lambda(G)\oplus\cdots\oplus\Lambda(G)}_{m\ \text{份}}
\subset \bigl(\RR^{E(G)}\bigr)^m,
\]
并以标准欧氏范数 \(|\cdot|_2\) 定义球截断
\[
\Lambda(G)^{\oplus m}(R):=\{h\in\Lambda(G)^{\oplus m}:\ |h|_2\le R\}.
\]
则当 \(R\to\infty\) 时有对数渐近式
\[
\boxed{\;
\log\bigl|\Lambda(G)^{\oplus m}(R)\bigr|
=mr\log R-\frac{m}{2}\log\tau(G)+O(1).\;}
\]
并且同一渐近式对任意仿射陪集 \(h_0+\Lambda(G)^{\oplus m}\) 的球截断计数同样成立。
\end{corollary}
\begin{proof}
由定理 \ref{thm:hypercube-cycle-lattice-theta-asymptotic}，\(\Lambda(G)\) 为秩 \(r\) 的欧氏格，且 \(\operatorname{covol}(\Lambda(G))=\sqrt{\tau(G)}\)。
直和格 \(\Lambda(G)^{\oplus m}\) 的秩为 \(mr\)，并且在正交直和的度量约定下满足协体积乘法性
\[
\operatorname{covol}\bigl(\Lambda(G)^{\oplus m}\bigr)=\operatorname{covol}(\Lambda(G))^{m}=\tau(G)^{m/2}.
\]
对格 \(\Lambda(G)^{\oplus m}\) 在其张成子空间中应用推论 \ref{cor:hypercube-cycle-lattice-counting-godel-log} 证明中的体积夹逼论证（将维数由 \(r\) 替换为 \(mr\)，协体积替换为 \(\tau(G)^{m/2}\)），得到
\[
\bigl|\Lambda(G)^{\oplus m}(R)\bigr|
=\frac{\operatorname{vol}(B_{mr})}{\tau(G)^{m/2}}\,R^{mr}+O(R^{mr-1}).
\]
取对数并吸收误差得所示对数渐近式。
对任意仿射陪集，平移不改变协体积与主导体积项，且截断边界层仅带来 \(O(R^{mr-1})\) 级别的修正，因此对数主项保持不变。证毕。
\end{proof}

\begin{corollary}[回路格截断的最小描述长度：树数对数扣除项]\label{cor:hypercube-cycle-lattice-directsum-mdl}
在推论 \ref{cor:hypercube-cycle-lattice-directsum-counting} 的设定下，取有限字母表 \(\mathcal A\)，\(\card{\mathcal A}=M\ge 2\)。
对任意 \(R\ge 2\)，若存在单射编码
\[
\mathrm{Enc}:\ \Lambda(G)^{\oplus m}(R)\hookrightarrow \mathcal A^{\le L},
\]
则必要长度满足
\[
L\ \ge\ \frac{\log\bigl|\Lambda(G)^{\oplus m}(R)\bigr|}{\log M}+O(1)
=\frac{mr\log R-\frac{m}{2}\log\tau(G)}{\log M}+O(1).
\]
\end{corollary}
\begin{proof}
单射性蕴含 \(\card{\mathcal A^{\le L}}\ge \bigl|\Lambda(G)^{\oplus m}(R)\bigr|\)。
另一方面，
\[
\card{\mathcal A^{\le L}}=\sum_{\ell=0}^{L}M^{\ell}=\frac{M^{L+1}-1}{M-1}\le \frac{M^{L+1}}{M-1},
\]
故 \(L\ge \frac{\log|\Lambda(G)^{\oplus m}(R)|}{\log M}+O(1)\)。
代入推论 \ref{cor:hypercube-cycle-lattice-directsum-counting} 的对数渐近式即得结论。证毕。
\end{proof}

\endinput

